\chapter{Paramétrer la surface}\label{ch:dv:parametrising-the-surface}

L'écran d'un teneur de marché affiche 400 strikes sur 20 échéances pour un même indice, chacun
avec un cours acheteur et un cours vendeur, plusieurs fois par seconde. Derrière l'écran, la
surface du desk compte cinq nombres par échéance. Elle ne passe pas par tous les
mids sont bruités, à un quart de point de volatilité près ici et là, et une courbe
qui passe par tous ces points se courbe entre eux de façon à créer
de l'arbitrage et faire varier les \omterm{def:dv:greeks-and-the-hedging-pnl:greeks}{grecques} de toutes les options évaluées à partir de celle-ci. Elle passe
par chaque fourchette, elle ne peut produire de densité négative, elle indique au
desk où vont les ailes au-delà du dernier strike coté, et ses cinq nombres peuvent
être surveillés~: lorsque le deuxième saute, c'est qu'il s'est passé quelque chose. Ce chapitre construit
cette surface. Il présente la paramétrisation employée dans toute la profession, sa
version à l'échelle de la surface, assortie de garanties contre l'arbitrage, le théorème qui borne
la pente que peuvent atteindre les ailes, la manière d'ajuster sur les cours acheteur et vendeur plutôt que
sur les mids, et les ajustements qu'exige une surface pour les dates de résultats et le week-end.

\section{Pourquoi paramétrer}

Le chapitre~7 construisait une surface par interpolation des cotations. Cela fonctionne lorsque les cotations sont
propres et denses. En pratique, elles ne le sont ni l'une ni l'autre~: les mids portent un bruit de l'ordre de la
fourchette, les strikes lointains sont cotés larges ou pas du tout, et chaque échéance a son
propre jeu de strikes. Un smile paramétrique fait quatre choses qu'une interpolation ne peut pas faire~:
\begin{itemize}
\item il lisse le bruit d'un strike à l'autre, au lieu de le reproduire~;
\item il peut être contraint à être exempt d'\omterm{def:dv:implied-volatility-and-its-surface:static}{arbitrage statique}~;
\item il extrapole au-delà des strikes cotés, de manière maîtrisée~;
\item il résume un smile en quelques nombres dont les variations d'un jour à l'autre peuvent être suivies
et expliquées.
\end{itemize}
Choisir les nombres pour que le smile colle aux cotations est une calibration, au sens du
livre~4, chapitre~24~: un petit problème d'optimisation, en général non
convexe.

\section{SVI et sa version de surface}

\begin{definition}[Paramétrisation SVI]\label{def:dv:parametrising-the-surface:svi}
La \emph{paramétrisation SVI}\index{paramétrisation SVI} (stochastic volatility
inspired) d'un smile écrit la \omterm{def:dv:implied-volatility-and-its-surface:surface}{variance implicite totale} d'une
\[
w(k)=a+b\Bigl(\rho(k-m)+\sqrt{(k-m)^2+\varsigma^2}\Bigr),
\]
avec $b\ge0$, $|\rho|<1$, $\varsigma>0$ et $a+b\varsigma\sqrt{1-\rho^2}\ge0$ (de sorte que le
minimum de $w$ est positif ou nul). Chaque paramètre a un rôle~: $a$ le niveau, $b$
l'angle entre les ailes, $\rho$ la rotation (le skew), $m$ la position horizontale du
minimum, $\varsigma$ la régularité du sommet.
\end{definition}

Pour $|k|\to\infty$, la variance totale de SVI est linéaire en $k$, avec les pentes $b(1+\rho)$ vers
la droite et $b(1-\rho)$ vers la gauche. Cette forme n'est pas un hasard~: le théorème suivant dit
qu'aucun smile sans arbitrage ne peut croître plus vite, et que SVI peut atteindre toute pente admissible.

\begin{theorem}[Formule des moments de Lee]\label{thm:dv:parametrising-the-surface:lee}
Soit $\beta_R=\limsup_{k\to\infty}w(k)/k$. Alors $\beta_R\in[0,2]$, et elle est déterminée par le
nombre de moments finis du sous-jacent~: avec $\tilde p=\sup\{p:\E^{\mathbb Q}S_T^{1+p}<\infty\}$,
\[
\beta_R=2-4\Bigl(\sqrt{\tilde p^2+\tilde p}-\tilde p\Bigr).
\]
L'aile gauche obéit à la même loi, avec $\tilde q=\sup\{q:\E^{\mathbb Q}S_T^{-q}<\infty\}$. Il s'agit de la
\emph{formule des moments de Lee}\index{formule des moments de Lee} (Lee moment formula).
\end{theorem}
\admitted

La borne 2 n'est atteinte que si le sous-jacent n'a aucun moment fini au-delà du premier~; une
lognormale, qui possède tous ses moments, a une pente nulle. Un ajustement SVI dont la pente d'aile $b(1\pm\rho)$ dépasse
2 présente un arbitrage dans son extrapolation, quoi qu'il fasse à l'intérieur de la plage cotée, et un ajustement
contraint à $b(1+|\rho|)\le2$ ne peut en produire aucun au-delà
(voir la \cref{fig:dv:parametrising-the-surface:wings}).

\begin{omfigure}
\begin{tikzpicture}
\begin{axis}[width=10.5cm,height=4.8cm,
  xlabel={log-moneyness $k$},ylabel={variance totale $w$},
  tick label style={font=\small},label style={font=\small},
  xmin=-2,xmax=2,ymin=0,ymax=0.6,grid=major,grid style={gray!25},
  legend columns=2,legend style={font=\footnotesize,at={(0.5,-0.3)},anchor=north,draw=none,fill=none,/tikz/every even column/.append style={column sep=8pt}},
  table/col sep=comma]
\addplot[omThm,very thick] table[x=k,y=w]{figdata/derivatives/08-parametrising-the-surface/wings.csv};
\addplot[gray,thick,dashed] table[x=k,y=lee]{figdata/derivatives/08-parametrising-the-surface/wings.csv};
\draw[omDef,thick] (axis cs:-0.40,0) -- (axis cs:-0.40,0.6);
\draw[omDef,thick] (axis cs:0.25,0) -- (axis cs:0.25,0.6);
\node[font=\footnotesize,omDef,anchor=east,align=right] at (axis cs:-0.45,0.48) {plage cotée\\entre les lignes};
\legend{tranche SVI à trois mois ajustée,borne de Lee $2|k|$}
\end{axis}
\end{tikzpicture}

\omcaption{La tranche SVI à trois mois du tutoriel, prolongée bien au-delà de ses strikes cotés
(entre les lignes verticales). Ses ailes sont droites, de pentes 0.062 à gauche et 0.014 à
droite, très en deçà de la borne de Lee $2|k|$ (en pointillés)~: l'extrapolation ne peut créer d'arbitrage.
Données~: le tutoriel.}\label{fig:dv:parametrising-the-surface:wings}
\end{omfigure}

Un bon ajustement tranche par tranche ne fait pas une bonne surface~: deux tranches ajustées indépendamment peuvent
se croiser, ce qui constitue un \omterm{def:dv:implied-volatility-and-its-surface:static}{arbitrage calendaire} (chapitre~7). La version de surface règle cela en liant les
tranches entre elles par la variance totale à la monnaie.

\begin{definition}[Surface SVI]\label{def:dv:parametrising-the-surface:ssvi}
La \emph{surface SVI}\index{surface SVI} (SSVI) écrit l'ensemble de la surface comme fonction de
\omterm{def:dv:implied-volatility-and-its-surface:surface}{log-moneyness} et la variance totale à la monnaie $\theta_T$ de chaque échéance,
\[
w(k,\theta_T)=\frac{\theta_T}2\Bigl(1+\rho\varphi(\theta_T)k+\sqrt{\bigl(\varphi(\theta_T)k+\rho\bigr)^2+1-\rho^2}\Bigr),
\]
avec un paramètre de type corrélation $\rho$ et une fonction de courbure $\varphi$, généralement la
loi puissance $\varphi(\theta)=\eta\theta^{-\gamma}$.
\end{definition}

\begin{theorem}[Conditions d'absence d'arbitrage pour la SSVI]\label{thm:dv:parametrising-the-surface:ssvi}
Une surface SSVI dont $\theta_T$ est non décroissant en $T$ est exempte d'\omterm{def:dv:implied-volatility-and-its-surface:static}{arbitrage calendaire} si
$0\le\partial_\theta(\theta\varphi(\theta))\le\frac{1}{\rho^2}\bigl(1+\sqrt{1-\rho^2}\bigr)\varphi(\theta)$~;
pour la loi puissance, cela vaut dès que $0<\gamma<1$. Elle est exempte d'\omterm{def:dv:implied-volatility-and-its-surface:static}{arbitrage butterfly} si, pour
chaque $\theta$, $\theta\varphi(\theta)(1+|\rho|)<4$ et $\theta\varphi(\theta)^2(1+|\rho|)\le4$.
\end{theorem}
\admitted

Le théorème et sa démonstration figurent chez Gatheral et Jacquier~; la première condition butterfly est également
nécessaire. Avec la loi puissance et $\gamma=\tfrac12$, le skew de volatilité à la monnaie vaut
$\rho\eta/(2\sqrt T)$~: la décroissance en $T^{-1/2}$ des skews actions, que le chapitre~12 reprendra.

\begin{example}[Une surface, cinq nombres]\label{ex:dv:parametrising-the-surface:ssvi}
Calibrée sur les cinq échéances du tutoriel (un mois à deux ans), la SSVI donne $\rho=-0.560$,
$\eta=1.095$, $\gamma=0.478$ et satisfait les deux conditions. Elle manque les mids de 0.09
points de volatilité à un an, de 0.76 à trois mois et de 2.0 à un mois
(\cref{fig:dv:parametrising-the-surface:ssvi})~: les smiles courts du marché sont plus asymétriques que
un seul $\rho$ ne peut le décrire. Les desks utilisent la SSVI comme backbone exempt d'arbitrage et ajoutent des corrections tranche par tranche
partout où les cotations l'exigent.
\end{example}

\begin{omfigure}
\begin{tikzpicture}
\begin{axis}[width=10.5cm,height=5cm,
  xlabel={log-moneyness $k$},ylabel={volatilité implicite (\%)},
  tick label style={font=\small},label style={font=\small},
  xmin=-0.4,xmax=0.25,ymin=8,ymax=44,grid=major,grid style={gray!25},
  legend columns=3,legend style={font=\footnotesize,at={(0.5,-0.3)},anchor=north,draw=none,fill=none,/tikz/every even column/.append style={column sep=8pt}},
  table/col sep=comma]
\addplot[omProp,only marks,mark=*,mark size=1.3pt] table[x=k,y=m1]{figdata/derivatives/08-parametrising-the-surface/ssvi.csv};
\addplot[omDef,only marks,mark=square*,mark size=1.3pt] table[x=k,y=m3]{figdata/derivatives/08-parametrising-the-surface/ssvi.csv};
\addplot[omThm,only marks,mark=triangle*,mark size=1.6pt] table[x=k,y=y1]{figdata/derivatives/08-parametrising-the-surface/ssvi.csv};
\addplot[omProp,thick] table[x=k,y=m1fit]{figdata/derivatives/08-parametrising-the-surface/ssvi.csv};
\addplot[omDef,thick] table[x=k,y=m3fit]{figdata/derivatives/08-parametrising-the-surface/ssvi.csv};
\addplot[omThm,thick] table[x=k,y=y1fit]{figdata/derivatives/08-parametrising-the-surface/ssvi.csv};
\legend{mids à 1 mois,mids à 3 mois,mids à 1 an}
\end{axis}
\end{tikzpicture}

\omcaption{Une surface SSVI (traits) calibrée simultanément sur cinq échéances, représentée pour trois d'entre elles
face aux mids (marqueurs). L'ajustement est proche à un an et manque le skew à un mois, plus marqué
qu'un seul $\rho$ ne le permet. Données~: le tutoriel.}\label{fig:dv:parametrising-the-surface:ssvi}
\end{omfigure}

\section{Splines et ajustements non paramétriques}

L'alternative à une formule est une courbe souple~: une spline cubique en $k$ passant par les cotations, ou une
spline de lissage qui arbitre entre la proximité aux cotations et la rugosité. Une spline passant par les
mids les reproduit exactement, bruit compris~; sur la tranche à trois mois du tutoriel, elle s'écarte de
0.17 points de volatilité du smile dont les cotations sont issues, deux fois plus que l'ajustement SVI
(0.08), et sa dérivée seconde, qui est la densité, hérite du bruit
(\cref{fig:dv:parametrising-the-surface:fit}). Une spline de lissage supprime l'essentiel de cet écart au prix
d'un paramètre (le poids de lissage) qui n'a aucun sens de marché~; et aucune de ces deux splines ne dit
rien sur les ailes au-delà du dernier strike, et ne garantit pas l'absence d'arbitrage sans
contraintes explicites. Les splines servent là où aucune formule ne convient (un smile à deux bosses
autour d'un événement binaire) et comme contrôle d'un ajustement paramétrique.

\section{Ajustement aux cours acheteur et vendeur}

Un mid n'est pas un prix auquel quiconque traite. Ce que dit le marché, c'est que la fair value de chaque
option se situe entre son cours acheteur et son cours vendeur, et tout smile contenu dans chaque fourchette y est compatible.

\begin{method}[Ajuster la SVI à une fourchette de cotations]\label{met:dv:parametrising-the-surface:fit}
Pour une échéance, avec des cotations $w^{\mathrm{bid}}_j\le w^{\mathrm{ask}}_j$ en $k_j$~:
\begin{enumerate}
\item Objectif~: l'écart quadratique dont le modèle sort de chaque fourchette, plus un petit poids sur
l'écart au mid, pour départager~; pondérer chaque strike par l'inverse du carré de sa largeur, de sorte
que les cotations serrées pèsent davantage.
\item Problème interne~: à $(m,\varsigma)$ fixé, avec $y=(k-m)/\varsigma$, la SVI est linéaire en
$(a,\ b\rho\varsigma,\ b\varsigma)$~; résoudre par moindres carrés pondérés et projeter sur le domaine
$b(1+|\rho|)\le2$, $|\rho|<1$, $a\ge0$ (la méthode quasi-explicite).
\item Problème externe~: minimiser sur les deux paramètres restants $(m,\varsigma)$ par une recherche sans dérivée
à partir de quelques points de départ.
\item Vérifier le résultat~: chaque point ajusté dans sa fourchette (ou signaler ceux qui n'y sont pas), les ailes dans
la borne de Lee, le facteur de densité $g$ du chapitre~7 positif ou nul sur une grille fine.
\end{enumerate}
\end{method}

Sur la tranche à trois mois du tutoriel, 27 strikes avec des fourchettes de largeur 0.4 point de
volatilité à la monnaie à 2 points à l'extrême gauche, l'ajustement tombe dans les 27 fourchettes, avec un écart maximal à un mid
de 0.18 points, et des paramètres $a=0.0019$, $b=0.038$, $\rho=-0.63$, $m=0.059$, $\varsigma=0.066$.

\begin{omfigure}
\begin{tikzpicture}
\begin{axis}[width=10.5cm,height=5cm,
  xlabel={log-moneyness $k$},ylabel={volatilité implicite à 3 mois (\%)},
  tick label style={font=\small},label style={font=\small},
  xmin=-0.41,xmax=0.26,ymin=10,ymax=37,grid=major,grid style={gray!25},
  legend columns=3,legend style={font=\footnotesize,at={(0.5,-0.3)},anchor=north,draw=none,fill=none,/tikz/every even column/.append style={column sep=8pt}},
  table/col sep=comma]
\addplot[gray,only marks,mark=-,mark size=3pt,error bars/.cd,y dir=both,y explicit]
  table[x=k,y=mid,y error expr=\thisrow{ask}-\thisrow{mid}]{figdata/derivatives/08-parametrising-the-surface/quotes.csv};
\addplot[omThm,very thick] table[x=k,y=svi]{figdata/derivatives/08-parametrising-the-surface/fits.csv};
\addplot[omDef,thin,dashed] table[x=k,y=spline]{figdata/derivatives/08-parametrising-the-surface/fits.csv};
\legend{fourchette bid--ask,SVI ajustée aux fourchettes,spline passant par les mids}
\end{axis}
\end{tikzpicture}

\omcaption{Une tranche à trois mois~: 27 fourchettes bid--ask, la courbe SVI ajustée aux fourchettes et une spline cubique
passant par les mids. À cette échelle, les deux courbes se confondent~; la spline porte le bruit des mids
(jusqu'à 0.17 point du vrai smile) là où la SVI reste à moins de 0.08 point. Données~: le
tutoriel.}\label{fig:dv:parametrising-the-surface:fit}
\end{omfigure}

\section{Événements, et temps ouvré face au temps calendaire}

Une surface suppose que la variance s'accumule régulièrement avec le temps. Deux faits le contredisent~: les événements connus,
qui ajoutent un bloc de variance à une date donnée, et les week-ends et jours fériés, pendant lesquels presque aucune ne s'accumule.

\begin{definition}[Variance d'événement]\label{def:dv:parametrising-the-surface:event}
La \emph{variance d'événement}\index{variance d'événement} d'une annonce programmée (une publication de résultats, une décision
de banque centrale, un résultat d'essai) est la variance du log-rendement provoqué par l'annonce
elle-même. Entre deux échéances encadrant l'événement, c'est l'excès de la variance totale sur ce que le
taux diffusif de l'échéance antérieure s'accumulerait~:
$w_{\mathrm{event}}=w(T_2)-\bigl(w(T_1)/T_1\bigr)T_2$.
\end{definition}

\begin{example}[Une semaine de résultats]\label{ex:dv:parametrising-the-surface:earnings}
Une action publie ses résultats après la clôture le jour 7. Les volatilités à la monnaie valent 32.0~\% pour l'échéance à 4
jours, 56.1~\% à 11 jours, 48.2~\% à 18 et 44.3~\% à 25. La \omterm{def:dv:parametrising-the-surface:event}{variance d'événement} vaut
$0.561^2\times11/365-0.32^2\times11/365=0.0064$~: un écart-type de 8.0~\% pour le mouvement lors de
l'annonce, un mouvement absolu attendu de $8.0\%\times\sqrt{2/\pi}=6.4\%$. Une fois celle-ci retirée, la
volatilité de l'échéance à 18 jours est de 32.0~\%~: toute la structure par terme au-delà de l'événement se résume à une seule
volatilité diffusive plus un saut (\cref{fig:dv:parametrising-the-surface:earnings}).
\end{example}

Interpoler linéairement la variance totale en temps calendaire entre les échéances à 4 et à 11 jours répartirait
le saut sur toute la semaine, et ferait coter une option à 7 jours (expirant avant l'annonce) à
près de 50~\% au lieu de 32~\%. Les études sur les options sur actions individuelles documentent ce schéma~: la volatilité implicite
monte à l'approche des résultats, chute brutalement après leur publication (Dubinsky et Johannes). Les surfaces
portent donc un calendrier d'événements, chacun avec sa propre variance, ajustés à partir de la structure par terme
et retirés avant toute interpolation en temps.

\begin{omfigure}
\begin{tikzpicture}
\begin{axis}[width=10.5cm,height=4.8cm,
  xlabel={jours avant l'échéance},ylabel={volatilité à la monnaie (\%)},
  tick label style={font=\small},label style={font=\small},
  xmin=0,xmax=27,ymin=25,ymax=62,grid=major,grid style={gray!25},
  legend columns=2,legend style={font=\footnotesize,at={(0.5,-0.3)},anchor=north,draw=none,fill=none,/tikz/every even column/.append style={column sep=8pt}},
  table/col sep=comma]
\addplot[omThm,very thick,mark=*,mark size=2pt] table[x=days,y=vol]{figdata/derivatives/08-parametrising-the-surface/earnings.csv};
\addplot[omDef,thick,dashed] table[x=days,y=base]{figdata/derivatives/08-parametrising-the-surface/earnings.csv};
\draw[gray,dotted,thick] (axis cs:7,25) -- (axis cs:7,62);
\node[font=\footnotesize,anchor=west] at (axis cs:7.3,28) {résultats};
\legend{cotées,hors événement~: 32~\%}
\end{axis}
\end{tikzpicture}

\omcaption{Structure par terme à la monnaie d'une action autour de la publication de ses résultats (pointillés). Les échéances postérieures
à la publication portent un bloc de variance unique, valant 8.0~\% d'écart-type~; son poids décroît en
$1/T$ avec l'échéance. Cotations illustratives.}\label{fig:dv:parametrising-the-surface:earnings}
\end{omfigure}

\begin{definition}[Temps ouvré]\label{def:dv:parametrising-the-surface:business}
Le \emph{temps ouvré}\index{temps ouvré} est une horloge qui compte le temps porteur de variance entre
deux dates plutôt que les jours calendaires~: les jours de cotation comptent pour un, les week-ends et les jours fériés pour une fraction (souvent
zéro), et les jours d'événement pour un plus leur \omterm{def:dv:parametrising-the-surface:event}{variance d'événement}. Les volatilités cotées sur une horloge de 252 jours de cotation
sont exprimées en temps ouvré.
\end{definition}

Une option d'une semaine à 20~\% sur une horloge de 252 jours (cinq jours de cotation) vaut autant qu'une volatilité de 20.34~\%
sur l'horloge calendaire de 365 jours, sur une durée de sept jours. La différence se manifeste dans le \omterm{def:dv:greeks-and-the-hedging-pnl:greeks}{theta}~: sur l'horloge calendaire, un report
du vendredi au lundi coûte deux septièmes de la valeur temps d'une semaine, sur l'horloge ouvrée presque
rien, la perte se produisant alors les jours de cotation. Un desk qui valorise en temps calendaire et
voit le marché se recoter chaque lundi utilise la mauvaise horloge.

\section{Tutoriel~: ajuster une tranche en direct}\label{tut:dv:parametrising-the-surface:tut}

\begin{tutorial}
\textbf{Objectif.} Ajuster SVI aux cours acheteur et vendeur sur une échéance, SSVI à une surface entière, et retirer un événement de résultats
d'une structure par terme. \textbf{État final~:} les quatre figures du chapitre et les chiffres des
\cref{ex:dv:parametrising-the-surface:ssvi,ex:dv:parametrising-the-surface:earnings}.
\begin{tutsteps}
\item \textbf{Le problème interne}~: pour $(m,\varsigma)$ fixés, SVI est linéaire en trois nombres~; les résoudre et
projeter sur le domaine admissible.
\omcode{code/firm/svi/firm_svi.py}{58}{69}{Ajustement intérieur quasi explicite de SVI, avec la borne de Lee sur les ailes.}\label{lst:dv:parametrising-the-surface:inner}
\item \textbf{L'événement}~: variance excédant le taux de l'échéance précédente, et le mouvement qu'elle implique.
\omcode{code/firm/svi/firm_svi.py}{130}{139}{Variance d'événement entre deux échéances, et mouvement implicite.}\label{lst:dv:parametrising-the-surface:event}
\item \textbf{Exécuter} \texttt{dv\_svi.fit\_slice()}, \texttt{ssvi\_fit()}, \texttt{earnings\_problem()}
et \texttt{fig\_svi.py}.
\end{tutsteps}
\textbf{Pour aller plus loin.} Doubler chaque largeur de fourchette cours acheteur--vendeur et observer la dérive des paramètres ajustés dans
la bande jour après jour~; puis donner à l'ajustement SSVI un $\rho$ distinct pour les deux premières échéances et
comparer l'erreur à un mois.
\end{tutorial}

\section{Construction~: l'ajusteur SVI}\label{bld:dv:parametrising-the-surface:svi}

\begin{build}
\bfield{Rôle} Transforme les cotations brutes en surface de la petite firme~: une tranche SVI par échéance, ajustée aux
bandes, vérifiée, avec SSVI en repli et extrapolation là où les cotations manquent.
\bfield{Interface} \texttt{svi(k, a, b, rho, m, s)}~; \texttt{svi\_check}~; \texttt{fit\_svi(k, w\_mid,
w\_bid, w\_ask, weights) -> (params, rmse)}~; \texttt{ssvi(k, theta, rho, eta, gamma)}~;
\texttt{ssvi\_check}~; \texttt{fit\_ssvi(slices)}~; \texttt{event\_variance}~; \texttt{implied\_move}~;
\texttt{nelder\_mead}.
\bfield{Règles} Ajuster en variance totale~; borne de Lee $b(1+|\rho|)\le2$ imposée dans l'ajustement~; signaler les
strikes laissés hors de leur bande~; ne jamais transmettre en aval une tranche dont le facteur de densité est négatif.
\bfield{Tests d'acceptation} \texttt{code/firm/svi/tests/}~: une tranche SVI exacte est retrouvée à
$10^{-5}$ près en variance totale~; la tranche ajustée du tutoriel se trouve dans toutes ses bandes sur plusieurs
graines~; les conditions de SSVI sont vérifiées pour la loi de puissance avec $\gamma<1$ et échouent lorsque $\theta\varphi$ est
trop grand~; la \omterm{def:dv:parametrising-the-surface:event}{variance d'événement} retrouve un saut implanté.
\bfield{Pour aller plus loin} Ajustement conjoint de tranches SVI avec pénalités d'\omterm{def:dv:implied-volatility-and-its-surface:static}{arbitrage calendaire} entre elles, et un
réajustement quotidien qui part des paramètres de la veille (chapitre 24).
\end{build}

\begin{omsources}
\item J.~Gatheral, «~A parsimonious arbitrage-free implied volatility parameterization with
application to the valuation of volatility derivatives~», présentation, Global Derivatives, Madrid
(2004).
\item J.~Gatheral et A.~Jacquier, «~Arbitrage-free SVI volatility surfaces~», \emph{Quantitative
Finance} 14 (2014) 59--71.
\item R.~W.~Lee, «~The moment formula for implied volatility at extreme strikes~», \emph{Mathematical
Finance} 14 (2004) 469--480.
\item C.~Martini et S.~De~Marco, «~Quasi-explicit calibration of Gatheral's SVI model~», Zeliade
livre blanc (2009, révisé en 2012).
\item A.~Dubinsky et M.~Johannes, «~Earnings announcements and equity options~», document de travail,
Columbia University (2006).
\end{omsources}

\section{Exercices}

\begin{exercise}[$\star$]\label{exo:dv:parametrising-the-surface:1}
Pour la tranche trimestrielle ajustée ($a=0.0019$, $b=0.038$, $\rho=-0.63$, $m=0.059$, $\varsigma=0.066$),
donner la variance totale et la volatilité implicite à la monnaie ($k=0$).
\end{exercise}

\begin{exercise}[$\star$]\label{exo:dv:parametrising-the-surface:2}
Donner les pentes des ailes gauche et droite de cette tranche, ainsi que le nombre de moments négatifs finis que
la pente gauche implique par la formule de Lee.
\end{exercise}

\begin{exercise}[$\star$]\label{exo:dv:parametrising-the-surface:3}
L'échéance précédant une annonce (5 jours) est à 30~\% et celle qui la suit (12 jours) à 50~\%. Donner la \omterm{def:dv:parametrising-the-surface:event}{variance
d'événement}, son écart-type et le mouvement absolu attendu.
\end{exercise}

\begin{exercise}[$\star\star$]\label{exo:dv:parametrising-the-surface:4}
Pourquoi est-il faux d'interpoler linéairement la variance totale en temps calendaire entre deux échéances encadrant
une publication de résultats~? Quel effet a-t-elle sur une option arrivant à échéance la veille de la publication~?
\end{exercise}

\begin{exercise}[$\star\star$]\label{exo:dv:parametrising-the-surface:5}
Avec les paramètres SSVI calibrés, vérifier les deux conditions d'absence d'\omterm{def:dv:implied-volatility-and-its-surface:static}{arbitrage butterfly} pour la variance
totale à la monnaie à un mois $\theta=0.00183$.
\end{exercise}

\begin{exercise}[$\star\star$]\label{exo:dv:parametrising-the-surface:6}
Une option à une semaine est cotée à 25~\% sur une horloge de 252 jours de négociation, avec cinq jours de
négociation restants. Donner la volatilité équivalente sur une horloge calendaire de 365 jours, sur une durée de sept jours.
\end{exercise}

\begin{exercise}[$\star\star\star$]\label{exo:dv:parametrising-the-surface:7}
\emph{Programmation.} Calibrer SVI sur les prix médians seulement, avec des poids égaux, et comparer le nombre de strikes situés
dans leur bande, ainsi que les pentes des ailes, avec celles de la calibration sur la bande.
\end{exercise}

\begin{exercise}[$\star\star\star$]\label{exo:dv:parametrising-the-surface:8}
\emph{Trouver le défaut.} «~La spline passe par chaque prix médian, donc son erreur d'ajustement est nulle~; SVI
s'écarte des prix médians d'au plus 0.18 point. La spline est la meilleure surface.~»
\end{exercise}

\section{Problème~: Semaine de publication des résultats}

\begin{problem}[{Problème du week-end --- le mouvement qu'intègrent les options}]\label{pb:dv:parametrising-the-surface:1}
Une action publie ses résultats après la clôture le jour 7. Ses volatilités implicites à la monnaie, cotées à 0.1 point
près, sont de 32.0~\% (4 jours), 56.1~\% (11 jours), 48.2~\% (18 jours) et 44.3~\% (25 jours). Un client demande
combien le marché anticipe que l'action va bouger lors de la publication, et ce que coûte un straddle acheté
pour cette publication.

\medskip\noindent\textbf{Partie I --- Variances totales.}
\begin{enumerate}
\item Donner la variance totale des échéances à 4 jours et à 11 jours.
\item Quelle variance totale aurait l'échéance à 11 jours en l'absence d'événement~?
\item Donner la \omterm{def:dv:parametrising-the-surface:event}{variance d'événement}.
\item Donner son écart-type et le mouvement absolu attendu d'un saut gaussien de cet écart-type.
\item Pourquoi utiliser l'échéance à 4 jours, et non celle à 25 jours, comme taux diffusif~?
\end{enumerate}

\medskip\noindent\textbf{Partie II --- Cohérence.}
\begin{enumerate}[resume]
\item Retirer l'événement de l'échéance à 18 jours~: quelle volatilité subsiste~?
\item Faire de même pour l'échéance à 25 jours.
\item Que suggérerait une volatilité résiduelle nettement supérieure à 32~\%~?
\item Que devient la volatilité à 11 jours le lendemain de la publication, si rien d'autre ne change~?
\item Pourquoi le straddle à 11 jours perd-il de la valeur pendant la nuit même si l'action bouge de 6~\%~?
\end{enumerate}

\medskip\noindent\textbf{Partie III --- La surface.}
\begin{enumerate}[resume]
\item Comment la surface doit-elle interpoler entre les échéances à 4 jours et à 11 jours~?
\item Les smiles de quelles échéances doivent être les plus pentus, et pourquoi~?
\item Comment traiter un week-end situé entre les échéances à 4 jours et à 11 jours~?
\item Une option à 7 jours arrive à échéance avant la publication~: quelle volatilité faut-il lui attribuer~?
\item Que deviennent les paramètres SVI calibrés de la tranche à 11 jours après la publication~?
\end{enumerate}

\medskip\noindent\textbf{Partie IV --- Jugement.}
\begin{enumerate}[resume]
\item Le 6.4~\% est-il la prévision du marché pour le mouvement~?
\item Comment comparer ce chiffre à l'historique des mouvements de l'action lors des publications de résultats~?
\item Quelle transaction traduit l'idée que le mouvement sera plus faible que ce qui est intégré dans les prix, et quel est son risque~?
\item Énoncer le \emph{résultat nommé}~: l'écart-type et la taille absolue attendue du mouvement
que les options intègrent dans leur prix.
\item En une phrase~: pourquoi faut-il retirer les événements avant d'interpoler dans le temps~?
\end{enumerate}
\end{problem}

\section{Questions d'entretien}

\begin{interviewq}[$\star$ \iqroles{chercheur, trader}]\label{iq:dv:parametrising-the-surface:1}
Écrivez SVI et dites ce que fait chaque paramètre au smile.
\end{interviewq}

\begin{interviewq}[$\star$ \iqroles{trader}]\label{iq:dv:parametrising-the-surface:2}
À partir des volatilités à la monnaie des échéances situées avant et après une publication de résultats, comment
obtenez-vous le mouvement implicite~?
\end{interviewq}

\begin{interviewq}[$\star\star$ \iqroles{chercheur}]\label{iq:dv:parametrising-the-surface:3}
Pourquoi la \omterm{def:dv:implied-volatility-and-its-surface:surface}{variance implicite totale} ne peut-elle pas croître plus vite que $2|k|$ dans les ailes~?
\end{interviewq}

\begin{interviewq}[$\star\star$ \iqroles{chercheur, développeur}]\label{iq:dv:parametrising-the-surface:4}
Vous calibrez SVI chaque seconde sur 20 échéances. Comment rendriez-vous la calibration rapide et stable, et
comment la calibreriez-vous sur les cours acheteur et vendeur plutôt que sur les prix médians~?
\end{interviewq}

\begin{interviewq}[$\star\star$ \iqroles{trader, risque}]\label{iq:dv:parametrising-the-surface:5}
Votre desk valorise le \omterm{def:dv:greeks-and-the-hedging-pnl:greeks}{theta} en temps calendaire. Que se passe-t-il chaque lundi matin, et que faudrait-il changer~?
\end{interviewq}

\begin{interviewq}[$\star\star\star$ \iqroles{chercheur}]\label{iq:dv:parametrising-the-surface:6}
Quand choisiriez-vous SSVI plutôt que des tranches SVI indépendantes, et à quoi renoncez-vous~?
\end{interviewq}
